\subsection{Relation to previous work}
\label{s-related-work}

\paragraph{Superabundant and colossally abundant numbers.}
Superabundant (SA) numbers are the image of the record map
\[
    X\mapsto\argmax_{n\leq X}\frac{\sigma(n)}n,
\]
and colossally abundant (CA) numbers are the image of the map
\[
    \eps\mapsto\argmax_{n\geq1}\frac{\sigma(n)}{n^{1+\eps}},
\]
defined for \(\eps>0\).
Ramanujan, Alaoglu and Erd\H{o}s, Nicolas, and Robin developed this
description and the associated benefit methods
\cite{ramanujan-hcn-1997,alaoglu1944,nicolas-1975-asterisque,
nicolas-1976-asterisque,robin-1980,robin-1983-rairo}.
Nicolas's marginal score \(F(p,k)\) governs whether the exponent of \(p\)
reaches \(k\) \cite[Section~4.1]{nicolas-2022-sigma}.

The score \(F(p,k)\) is also the value-to-log-size ratio in our knapsack
ordering. Fixing \(\omega(n)=m\) makes the first-exponent increments
\(F(p,1)\) mandatory for \(p\leq p_m\), while the optional increments
with \(k\geq2\) retain their classical order. The resulting optimizer is
CA and SA relative to its prescribed support, but need not be a global
CA or SA number. The bounded-support problem restores the support
decisions.

\paragraph{Classical Robin asymptotics.}
Under the Riemann hypothesis, Ramanujan's expansion for generalized
superior highly composite numbers, in the corrected version published by
Nicolas and Robin, already contains the normalized constant
\(2\sqrt2-2\) and the zeta-zero term that appear in the classical Robin
margin
\cite[formula~(382) and Note~71]{ramanujan-hcn-1997}. Nicolas gives a
modern effective treatment for CA numbers and then for all integers
\cite[Proposition~5.1 and Theorem~1.1]{nicolas-2022-sigma}. The
coefficient \(2\sqrt2-2\) and the zero sum in the final integer margin are
therefore not new here.

These results establish the classical integer margin, not the
fixed-support relaxation or its integrality gap. Our optimizer
characterization and rounding formula give that unconditional
gap. Nicolas's analytic estimates enter only afterward, when the gap is
combined with the relaxed value.

\paragraph{The Axler--Nicolas comparison.}
Axler and Nicolas \cite{axler-nicolas-2023} prove
\[
    \frac{\displaystyle\max_{n\leq X}n/\varphi(n)}
         {\displaystyle\max_{n\leq X}\sigma(n)/n}
    =
    1+
    \frac{2\sqrt2}
         {\sqrt{\log X}\log\log X}
    +
    O\left(
        \frac1{\sqrt{\log X}(\log\log X)^2}
    \right),
\]
where \(\varphi\) is Euler's totient function. Nicolas obtains a full
expansion \cite[Theorem~1.2]{nicolas-2025}.
Their size-constrained problem and our fixed-support penalized problem have
different feasible sets, so neither theorem implies the other. The
shared second-exponent score equation explains the same leading
coefficient. The second terms have different meanings in the two
problems. Section~\ref{s-score-connection} gives the precise comparison.

\paragraph{Other reductions.}
The first-prime support exchange is standard in work on possible
counterexamples to Robin's inequality
\cite[Section~5, Lemma~9]{choie2007}. Other work restricts a least
counterexample or reduces Robin's criterion to record subsequences
\cite{akbary-friggstad-2009,nazardonyavi-yakubovich-2014}. We instead
maximize separately on each level set of \(\omega(n)\). This makes the
first-exponent increments for the first \(m\) primes mandatory and leaves
the exponent lattice as the discrete part of the problem.

\paragraph{Outline.}
Section~\ref{s-fixed-slice} formulates the problem.
Section~\ref{s-optimizers} gives the scalar relaxed optimizer and the
score-ordered prefix theorem.
Section~\ref{s-integrality-gap} proves one-step proximity and evaluates
the rounding loss. Its final subsection relates the calculation
to CA shells and the Axler--Nicolas comparison.
Sections~\ref{s-relaxed-excess} and \ref{s-robin-margin} add Nicolas's
analytic input and derive the fixed-slice criterion.
Section~\ref{s-bounded-support} treats \(\omega(n)\leq m\).
